\begin{table}[!htbp]
\centering
\caption{Weighting and dependence sensitivity}\label{tab:dependence}
\small
\setlength{\tabcolsep}{5pt}
\begin{tabular}{llrrlrrr}
\toprule
Type & Period & Event & Review & Leave one out & CR1 & CV3 & Two way \\
\midrule
Addition & 2017--2026 & 2.15 & 2.50 & [1.87, 2.47] & 0.55 & 0.55 & 0.47 \\
Addition & 2022--2026 & 1.50 & 1.50 & [1.02, 2.01] & 0.74 & 0.77 & 0.71 \\
Deletion & 2017--2026 & -1.04 & -0.52 & [-1.45, -0.50] & 0.87 & 0.88 & 0.87 \\
Deletion & 2022--2026 & -1.53 & -1.32 & [-2.06, -0.58] & 1.12 & 1.17 & 1.11 \\
\bottomrule
\end{tabular}
\begin{minipage}{\linewidth}\footnotesize\vspace{5pt}
CAR $[AD-1,ED+1]$ in percent. Review mean gives each review equal weight. Leave one out is the range of event-weighted estimates after removing one review at a time, not a confidence interval. CR1, CV3 and Two way report standard errors. CV3 is the cluster jackknife centered on the full-sample estimate. Two way combines review and security clustering with subtraction of the review-security intersection. It does not solve dependence between different securities across different reviews. Source: audit/results/dependence\_sensitivity.csv.
\end{minipage}
\end{table}
